\documentclass[10pt,a4paper]{amsart}
\usepackage[margin=19mm]{geometry}
\usepackage{amsmath,amssymb,mathtools}
\usepackage[scr=rsfs,cal=euler]{mathalfa}
\usepackage{xfrac}
\usepackage[unicode,hidelinks,bookmarks=false]{hyperref}
\newcommand{\ie}{i.e.\ }
\newcommand{\cf}{cf.\ }
\newcommand{\resp}{respectively\ }
\newcommand{\Subsection}{\subsection}
\providecommand{\nobreakdash}{\nobreak\hbox{-}}
\setlength{\emergencystretch}{2em}
\multlinegap0pt
\usepackage{bm}
\usepackage{bbm}
\usepackage{dsfont}
\usepackage{xspace}
\usepackage{cancel}
\usepackage{mathtools}
\usepackage[shortlabels]{enumitem}
\usepackage{amsthm}
\makeatletter
\@ifundefined{theorem}{\newtheorem{theorem}{Theorem}}{}
\@ifundefined{proposition}{\newtheorem{proposition}[theorem]{Proposition}}{}
\makeatother
\newcommand{\NN}{\mathbb{N}}
\newcommand{\cl}{\overline}
\newcommand{\diam}{\textup{diam}}
\title[Questions in linear recurrence: From the $T\oplus T$-problem]{Questions in linear recurrence: From the $T\oplus T$-problem to lineability}
\author{Sophie Grivaux, Antoni L\'opez-Mart\'inez, Alfred Peris}
\date{Source: arXiv:2212.03652v3 (CC BY 4.0)}
\begin{document}
\maketitle

\noindent\emph{Source and licence.} Questions in linear recurrence: From the $T\oplus T$-problem to lineability. Version arXiv:2212.03652v3 (CC BY 4.0), distributed under the Creative Commons Attribution 4.0 International licence, as recorded in the arXiv licence metadata for this exact version and shown on the abstract page https://arxiv.org/abs/2212.03652v3. Full rights notice: licenses/arxiv-2212.03652-CC-BY-4.0.txt.\par\smallskip

\noindent\textit{Editorial scope.} Counted unit: the Proposition whose source label is \texttt{Pro:qr.complete} in the article (source0.tex lines 273--280, source label \texttt{Pro:qr.complete}), delivered together with the article's own complete original proof of it (source0.tex lines 281--317, one \texttt{proof} environment of 37 lines). No mathematics is written, repaired or re-derived: statement and proof are the article's own text line for line. Required context retained uncounted: none. Every definition, display and label of those lines is retained unchanged. Omitted from this package and not redistributed: 1--272, 318--1733. Nothing inside a retained excerpt is omitted or altered: every retained body is delivered exactly as the article's lines are written, so no omission receipt of any kind accompanies this package. Citations: the retained excerpts carry no citation command at all, so no bibliography entry of the article is transcribed and no reference is redistributed. Reference adaptation: none was needed and none was made; every reference inside the delivered unit resolves inside it, so no unresolved reference or citation is produced. Printed slips kept literal: no printed word, symbol, hypothesis or number of the counted statement or of its proof was altered. Layout and class: the article's own class is \texttt{article} with packages including graphicx, inputenc, babel, here, amsmath, amssymb, amscd, amsthm, mathrsfs, enumerate, hyperref, footmisc; the wrapper is the lane's pinned \texttt{amsart} editorial wrapper at 10pt on A4, which restates the theorem-like environments the retained text needs and adds the article's own macro definitions that the retained text needs (\texttt{\textbackslash NN}, \texttt{\textbackslash cl}, \texttt{\textbackslash diam}), with extra wrapper packages (bm, bbm, dsfont, xspace, cancel, mathtools, enumitem). The delivered numbering is the unit's own. Assets: no retained span contains a figure environment, a graphics-inclusion command, a PDF-page inclusion, a table or a TikZ picture; no image file is copied into this package, the package declares no asset dependency and no image file is redistributed with it. Licence: version arXiv:2212.03652v3 of this article is distributed under the Creative Commons Attribution 4.0 International licence, as recorded in the arXiv licence metadata for this exact version and shown on the abstract page https://arxiv.org/abs/2212.03652v3. A full rights notice is delivered at licenses/arxiv-2212.03652-CC-BY-4.0.txt. Characters: every retained excerpt is pure ASCII, so the delivered engine, XeLaTeX with its default Latin Modern text font, reports no missing character.
\begin{proposition}\label{Pro:qr.complete}
	Let $(X,T)$ be a Polish dynamical system. The following are equivalent:
	\begin{enumerate}[{\em(i)}]
		\item $T$ is quasi-rigid;
		
		\item $T$ is topologically quasi-rigid.
	\end{enumerate}  
\end{proposition}

\begin{proof}
	The implication (i) $\Rightarrow$ (ii) is straightforward: assume that there is a dense subset $Y \subset X$ and an increasing sequence $(n_k)_{k\in\NN}$ such that $T^{n_k}x \to x$ for every $x \in Y$. For each non-empty open subset $U \subset X$ we can select $x \in U \cap Y$, and then
	\[
	T^{n_k}x \in T^{n_k}(U) \cap U \quad \text{ for every sufficiently large } k \in \NN.
	\]
	Hence $T$ is topologically quasi-rigid with respect to $(n_k)_{k\in\NN}$. Let us prove (ii) $\Rightarrow$ (i): assume that $d(\cdot,\cdot)$ is a metric defining the complete topology of $X$, let $(U_s)_{s\in\NN}$ be a countable base of open sets for $X$ and set $V_{s,s}:=U_s$ for each $s \in \NN$. By (ii) we can find $n_1 \in \NN$ such that
	\[
	T^{n_1}(U_1) \cap U_1 = T^{n_1}(V_{1,1}) \cap V_{1,1} \neq \varnothing.
	\]
	The continuity of $T$ ensures the existence of a non-empty open subset $V_{1,2} \subset X$ of diameter (with respect to $d$) less than $\frac{1}{2^2}$ for which $\cl{V_{1,2}} \subset V_{1,1}$ and $T^{n_1}\left(\cl{V_{1,2}}\right) \subset V_{1,1}$. Suppose now that for some $k \in \NN$ we have already constructed:
	\begin{enumerate}[--]
		\item finite sequences $(V_{s,j})_{s\leq j\leq k}$ of open subsets of $X$, for each $1\leq s\leq k-1$;
		
		\item and a finite increasing sequence of positive integers $(n_j)_{1\leq j\leq k-1}$;
	\end{enumerate}
	with the properties that $V_{s,j}$ has $d$-diameter less than $\frac{1}{2^j}$ when $s<j$, and also that
	\[
	\cl{V_{s,j}} \subset V_{s,j-1} \quad \text{ and } \quad T^{n_{j-1}}\left(\cl{V_{s,j}}\right) \subset V_{s,j-1} \quad \text{ for all } 1\leq s < j \leq k.
	\]
	Then, considering the open subsets $V_{1,k}, V_{2,k}, ..., V_{k,k} \subset X$ and using again (ii) we can select a positive integer $n_k \in \NN$ with $n_k>n_{k-1}$ for which $T^{n_k}(V_{s,k}) \cap V_{s,k} \neq \varnothing$ for all $1\leq s\leq k$. Again the continuity ot $T$ ensures the existence, for each $1\leq s\leq k$, of a non-empty open subset $V_{s,k+1} \subset X$ of $d$-diameter less than $\frac{1}{2^{k+1}}$ and such that
	\[
	\cl{V_{s,k+1}} \subset V_{s,k} \quad \text{ and } \quad T^{n_{k}}\left(\cl{V_{s,k+1}}\right) \subset V_{s,k} \quad \text{ for all } 1\leq s\leq k.
	\]
	A recursive argument gives us an increasing sequence of positive integers $(n_k)_{k\in\NN}$ and, for each $s \in \NN$, a sequence of non-empty open sets $(V_{s,k})_{k=s}^{\infty}$ such that each set $V_{s,k}$ has $d$-diameter less than $\frac{1}{2^k}$ when $s<k$, but also satisfying
	\[
	\cl{V_{s,k+1}}\subset V_{s,k} \quad \text{ and } \quad  T^{n_k}\left(\cl{V_{s,k+1}}\right)\subset V_{s,k} \quad \text{ for all } s,k \in \NN \text{ with } s\leq k.
	\]
	By the Cantor intersection theorem, for each $s\in\NN$ there is a unique vector $y_s\in X$ with
	\[
	\{y_s\} = \bigcap_{k\geq s} \cl{V_{s,k}} \subset V_{s,s} = U_s.
	\]
	We deduce that $Y:=\{ y_s: s \in \NN \}$ is a dense subset of $X$. Since for each $s \in \NN$ we have that $T^{n_k}y_s \in T^{n_k}\left(\cl{V_{s,k+1}}\right) \subset V_{s,k}$ for all $k > s$, we get that
	\[
	\limsup_{k\to\infty} d(T^{n_k}y_s,y_s) \leq \limsup_{k\to\infty} \left(\diam_d(V_{s,k})\right) \leq \limsup_{k\to\infty} \frac{1}{2^k} = 0,
	\]
	and hence $T$ is quasi-rigid with respect to the sequence $(n_k)_{k\in\NN}$.
\end{proof}

\end{document}
